\FloatBarrier
\subsection{3b) Pong-v5 With a Baseline}

\textbf{Setup:} everything is identical to Part~1: network, frame difference, $\gamma = 0.99$, return restart after every point, learning rate 0.001, 5 episodes per update, entropy bonus, clipping, 2,000 episodes and seed 0. The only change is the weight of the log-probabilities in the actor's loss. Pong training is too slow for several seeds, so the comparison is \textbf{one run per variant} (see the limitation in the summary).

\subsubsection{Chosen Baseline and Implementation}

\textbf{Non-constant baseline: a learned state value $V(s_t)$ (actor-critic).} The actor is weighted with the advantage
\[
  A_t = G_t - V(s_t),
\]
where $V$ is a second network, the \emph{critic}, which learns to predict the return of the current rally.

\begin{center}
\small
\begin{tabularx}{\textwidth}{@{}p{4.2cm}L@{}}
  \toprule
  \textbf{Design choice} & \textbf{Reason} \\
  \midrule
  Critic: same layers as the policy, own Adam optimizer (lr 0.001, clipped to norm 1) & Trained on the same Monte-Carlo returns with the loss $\frac{1}{T}\sum_t \big(V(s_t) - G_t\big)^2$, updated together with the actor every 5 episodes. \\
  Monte-Carlo return $G_t$, not a bootstrapped TD target & $V(s_t)$ does not depend on the action $a_t$, so $\E\big[\nabla \log \pi(a_t \mid s_t)\, V(s_t)\big] = 0$: the expected gradient is the same as in Part~1, only its variance changes. A TD target would also reduce the variance, but bias the gradient. \\
  Critic input: current frame \textbf{and} frame difference ($2 \times 80 \times 80$) & The value of a state depends on where the paddle is relative to the ball, which the frame difference alone mostly hides (Part~1, Test 4). Any input that does not depend on $a_t$ keeps the baseline unbiased. \\
  $V$ is detached in the advantage & The actor's loss does not train the critic, and the critic's loss does not change the actor. \\
  State-dependent instead of constant & Returns vary strongly within an episode: close to $-1$ just before a lost point, close to $+1$ before a won one. A constant cannot follow this; $V(s_t)$ can. \\
  \bottomrule
\end{tabularx}
\end{center}

\textbf{Implementation:} episodes are played exactly as in Part~1; the critic's inputs are additionally stored as 8-bit integers. After each episode, the critic computes $V(s_t)$ for all steps in chunks of 256 steps and accumulates its gradient. The actor's loss is then
\[
  -\sum_t \log \pi(a_t \mid s_t)\,\big(G_t - V(s_t)\big) \;-\; 0.01 \sum_t \mathcal{H}_t ,
\]
and every 5 episodes both networks are clipped and updated.

\subsubsection{Training and Evaluation}

\fig{1.0}{3b_baseline_training.png}{Pong training without (left) and with (right) the learned baseline: episode reward and its moving average over 100 episodes.}{fig:3b-bl-train}

\fig{1.0}{3b_baseline_comparison.png}{Moving averages over 100 episodes of the reward, the episode length and the standard deviation of $P(\text{down})$, without and with the baseline.}{fig:3b-bl-comp}

\begin{itemize}
  \item \textbf{Learning starts much earlier.} With the baseline, the policy starts reacting to the screen (std of $P(\text{down})$) at episode $\sim$250 instead of $\sim$750, and the reward rises from the start. The moving average reaches $-19.5$ at episode \textbf{990 instead of 1,530} (Figure~\ref{fig:3b-bl-comp}).
  \item \textbf{It learns faster for the whole run.} The last 100 episodes reach \textbf{$-14.93$ instead of $-18.07$}, and the rallies are 50\% longer (3,028 vs.\ 2,065 steps).
  \item \textbf{Still improving at the end,} even more steeply than without baseline.
\end{itemize}

\fig{1.0}{3b_baseline_histogram.png}{Evaluation histograms without and with the baseline (500 episodes with sampled actions each).}{fig:3b-bl-hist}

\begin{center}
\small
\begin{tabular}{@{}lrr@{}}
  \toprule
  & \textbf{Without baseline} & \textbf{With baseline (actor-critic)} \\
  \midrule
  Mean training reward & $-19.83$ & $\mathbf{-18.74}$ \\
  Reward, last 100 training episodes & $-18.07$ & $\mathbf{-14.93}$ \\
  Moving average reaches $-19.5$ at episode & 1,530 & \textbf{990} \\
  \textbf{Evaluation, sampled (500 episodes)} & $\mathbf{-17.66 \pm 1.86}$ & $\mathbf{-13.27 \pm 2.87}$ \\
  Evaluation, greedy (50 episodes) & $-15.68 \pm 2.37$ & $\mathbf{-10.22 \pm 3.33}$ \\
  Greedy: points won per game / length & 5.3 / 2,960 & \textbf{10.8 / 3,847} \\
  \emph{Reference:} random / ball tracker (Part~1, Test 3) & $-20.40$ / $-14.74$ & \\
  \bottomrule
\end{tabular}
\end{center}

\begin{itemize}
  \item \textbf{The evaluation improves by 4.4 points} (Figure~\ref{fig:3b-bl-hist}). Every one of the 500 episodes is better than the random policy, and the best one reaches $-1$. The policy still loses every game.
  \item \textbf{It now beats the hand-coded ball tracker} ($-13.27$ sampled, $-10.22$ greedy, vs.\ $-14.74$). Without the baseline, it stayed clearly below it.
  \item \textbf{The larger std (2.87 vs.\ 1.86)} mostly reflects the wider range of results: a better player has more room between losing 21:0 and winning many points.
\end{itemize}

\subsubsection{What the Critic Learns}

\fig{1.0}{3b_baseline_critic.png}{The critic during training: its loss, its explained variance, and the mean absolute weight of the actor's update with and without baseline.}{fig:3b-bl-critic}

\begin{itemize}
  \item \textbf{First, the critic learns the average return.} Its loss drops from 0.8 to $\sim$0.12 within $\sim$150 episodes (Figure~\ref{fig:3b-bl-critic}). Most returns are negative (the agent loses most points), and subtracting this average alone shrinks the actor's weights from a mean $|G_t|$ of $\sim$0.78 to a mean $|A_t|$ of $\sim$0.22 within the first 200 episodes.
  \item \textbf{Then it learns to tell states apart.} The explained variance $1 - \Var(G - V)/\Var(G)$ stays below 0.05 for the first 1,000 episodes and rises to \textbf{0.26} over the last 100. The critic increasingly predicts \emph{which} rallies will be won.
  \item \textbf{The critic's loss rises slowly again} ($0.12 \to 0.21$). As the agent wins more points, the returns spread over the whole range from $-1$ to $+1$ and become harder to predict.
\end{itemize}

\fig{1.0}{3b_baseline_policy_response.png}{Direction gap (left) and position gap (right) over training, without and with the baseline.}{fig:3b-bl-response}

\begin{center}
\small
\begin{tabular}{@{}rrrrr@{}}
  \toprule
  & \multicolumn{2}{c}{\textbf{Direction gap}} & \multicolumn{2}{c}{\textbf{Position gap}} \\
  \cmidrule(lr){2-3} \cmidrule(l){4-5}
  \textbf{Checkpoint} & \textbf{Without} & \textbf{With baseline} & \textbf{Without} & \textbf{With baseline} \\
  \midrule
  250 & 0.000 & $\mathbf{+0.015}$ & 0.000 & $+0.006$ \\
  500 & $+0.003$ & $\mathbf{+0.095}$ & $+0.001$ & $-0.008$ \\
  1,000 & $+0.095$ & $+0.064$ & $-0.027$ & $-0.009$ \\
  1,500 & $+0.089$ & $+0.057$ & $-0.073$ & $+0.001$ \\
  2,000 & $+0.085$ & $+0.060$ & $-0.057$ & $\mathbf{+0.044}$ \\
  \bottomrule
\end{tabular}
\end{center}

\begin{itemize}
  \item \textbf{The direction gap appears about 500 episodes earlier} (+0.095 at episode 500, a level Part~1 only reaches at episode 1,000; Figure~\ref{fig:3b-bl-response}). It is the same feature as in Part~1, learned faster.
  \item \textbf{The position gap becomes positive at the end} (+0.044 at episode 2,000, vs.\ $-0.057$ without baseline). This is a first sign that the policy also starts moving \emph{towards} the ball, but it is below the threshold of 0.05 used in Part~1 and based on 8 episodes, so it is only a hint.
\end{itemize}

\subsubsection{Edge Cases}

The edge cases of Part~1 (sparse rewards, frame difference, saturation, end of an episode, sticky actions) apply unchanged.

\begin{center}
\small
\begin{tabularx}{\textwidth}{@{}p{3.6cm}L@{}}
  \toprule
  \textbf{Edge case} & \textbf{Handling} \\
  \midrule
  Biased baseline & A baseline must not depend on the action $a_t$. The critic only sees the frames \emph{before} $a_t$, and its value is detached in the advantage (tested). The Monte-Carlo target keeps the gradient unbiased. \\
  Untrained critic at the start & The untrained critic predicts $\sim$+0.09, while the returns average $\sim$$-0.76$. For the first $\sim$100 episodes it is therefore a bad baseline: the mean $|A_t|$ is even slightly larger than $|G_t|$, and the gradient variance is 1.2$\times$ that without baseline (Test 8). It learns the average within $\sim$150 episodes. I did not pre-train the critic, so both runs start identically. \\
  Critic lags behind the improving policy & The critic is trained on returns of earlier, worse policies, so it is too pessimistic (mean $V$ $-0.44$ vs.\ mean $G$ $-0.25$, Test 7). This does not bias the gradient; it only makes the variance reduction smaller than with a perfect critic. \\
  Unbounded critic output & $G_t \in [-1, 1]$, but the critic's output is not bounded. 1\% of its predictions lie outside $[-1, 1]$ (lowest $-1.41$). As a baseline this is harmless, so I did not add a $\tanh$. \\
  Long episodes (up to $\sim$4,000 steps) & The critic's forward and backward pass runs in chunks of 256 steps, and its inputs are stored as 8-bit integers (12.8\,KB per step). The chunked pass gives exactly the same values, loss and gradients as one pass (tested). \\
  Fair comparison & Same seed and same initial actor weights. The critic is initialised without consuming the random numbers used for sampling, so the first 5 episodes (before the first update) are identical to Part~1 (tested). \\
  Wrong optimizer & A critic optimizer that does not optimise exactly the critic's parameters would silently leave the critic untrained. It raises an error (tested), as do the invalid inputs of Part~1. \\
  Reproducibility & Training twice with the same seed gives identical results, including the critic's loss (tested). \\
  \bottomrule
\end{tabularx}
\end{center}

\subsubsection{Tests for Validating Correctness}

The tests continue the numbering of Part~1, and \textbf{all of them pass}.

\textbf{6. Fair setup and advantage computation.}
\begin{itemize}
  \item The actor starts from exactly the same weights as in Part~1, and the first 5 episodes are identical (rewards $-20, -21, -21, -21, -21$) \pass. Every later difference is caused by the baseline.
  \item The values in the advantage carry no gradient \pass.
  \item The chunked critic pass equals a single pass (values, loss and gradients) \pass.
  \item The explained variance is 1 for $V = G$ and 0 for $V = \bar G$ \pass.
\end{itemize}

\textbf{7. Does the critic predict the return?} 10 episodes of the final actor, evaluated with the untrained and the final critic on the same steps (Figure~\ref{fig:3b-bl-values}):

\begin{center}
\small
\begin{tabular}{@{}lrr@{}}
  \toprule
  & \textbf{Untrained critic} & \textbf{Final critic} \\
  \midrule
  Explained variance & 0.000 & \textbf{0.325} \pass \\
  Mean $V$ over the 20 steps before a \textbf{won} point (91 points) & & \textbf{+0.10} \\
  Mean $V$ over the 20 steps before a \textbf{lost} point (210 points) & & $\mathbf{-0.61}$ \pass \\
  \bottomrule
\end{tabular}
\end{center}

\fig{1.0}{3b_baseline_values.png}{Return and critic prediction in a test episode (left), and calibration of the final critic (right).}{fig:3b-bl-values}

\begin{itemize}
  \item The final critic clearly separates rallies that will be won from those that will be lost, and in the test episode its value rises before every won point.
  \item \textbf{Calibration:} the actual return increases monotonically with the predicted value (all 10 deciles), but lies 0.2--0.4 above it. The critic is too pessimistic, because it lags behind the improving policy (see the edge cases).
\end{itemize}

\textbf{8. Variance of the gradient estimate.} As for CartPole: for a fixed actor and critic, 20 batches of 5 episodes, with the gradient of every batch computed on the same episodes without baseline, with a constant baseline (the batch's mean return) and with the critic. The entropy term is left out.

\fig{0.7}{3b_baseline_gradient_variance.png}{Total variance of the gradient estimate for the initial and the final actor-critic checkpoint.}{fig:3b-bl-var}

\begin{center}
\small
\begin{tabular}{@{}llrrr@{}}
  \toprule
  \thead[l]{Checkpoint} & \thead[l]{Baseline} & \thead[r]{Total\\variance} & \thead[r]{Relative\\to none} & \thead[r]{Cosine of mean\\gradient with none} \\
  \midrule
  After 0 episodes & None & $3.5 \cdot 10^3$ & 1.00 & 1.00 \\
  & Constant (batch mean) & $5.3 \cdot 10^2$ & 0.15 & 0.70 \\
  & Critic (untrained) & $4.2 \cdot 10^3$ & 1.20 & 1.00 \\
  \addlinespace
  After 2,000 episodes & None & $4.7 \cdot 10^5$ & 1.00 & 1.00 \\
  & Constant (batch mean) & $2.5 \cdot 10^5$ & 0.52 & 0.99 \\
  & \textbf{Critic (trained)} & $\mathbf{1.4 \cdot 10^5}$ & \textbf{0.29} \pass & 0.96 \\
  \bottomrule
\end{tabular}
\end{center}

\begin{itemize}
  \item \textbf{The trained critic reduces the variance by a factor of 3.5}, clearly more than the constant baseline (factor 1.9) \pass. Once the agent wins points, a constant cannot follow the returns, which now vary between rallies.
  \item \textbf{At the start, the constant baseline is better and the untrained critic slightly harmful.} Almost all returns are negative, so their mean removes most of the noise; the untrained critic subtracts the wrong constant (+0.09). This is the same effect as in the edge cases. During training, the critic learns this mean within the first $\sim$150 episodes.
  \item \textbf{Same direction:} the mean gradients with the baselines point in the same direction as without (cosine 0.96--1.00), as expected for unbiased baselines. With 20 batches this is only a rough check; the 0.70 of the constant baseline at the start comes from the noise of the untrained policy's small gradient.
  \item The variance grows $\sim$130$\times$ over training, because episodes are $\sim$3$\times$ longer and the policy's gradients larger (Part~1, Figure~\ref{fig:3b-diag}).
\end{itemize}

\FloatBarrier
\begin{summarybox}{Summary of 3b, Part 2}
\begin{tabularx}{\linewidth}{@{}Lrr@{}}
  \toprule
  & \textbf{Without baseline} & \textbf{With baseline (actor-critic)} \\
  \midrule
  Policy starts reacting to the screen & episode $\sim$750 & \textbf{episode $\sim$250} \\
  Moving average reaches $-19.5$ & episode 1,530 & \textbf{episode 990} \\
  Reward, last 100 training episodes & $-18.07$ & $\mathbf{-14.93}$ \\
  Evaluation, sampled (500 episodes) & $-17.66 \pm 1.86$ & $\mathbf{-13.27 \pm 2.87}$ \\
  Evaluation, greedy (50 episodes) & $-15.68 \pm 2.37$ & $\mathbf{-10.22 \pm 3.33}$ \\
  Gradient variance, final policy & 1 & \textbf{0.29} \\
  \bottomrule
\end{tabularx}

\medskip
The learned baseline makes REINFORCE on Pong learn much faster: the policy starts reacting to the ball about 500 episodes earlier, and after the same 2,000 episodes it wins about twice as many points and beats the hand-coded ball tracker. The critic learns first the average return and later which rallies will be won, and in the end reduces the gradient variance by a factor of 3.5.

\textbf{Limitation:} this is a single run per variant. The difference is large and consistent over the whole run, but without several seeds I cannot quantify how much of it depends on the seed.
\end{summarybox}
